%%%%%%%%%%%%%%%%%%%%%%%%%%%%%%%%%%%%%%%%%%%%%%%%%%%%%%%%%%%%%%%%%%%%%%%%%%%%%%%
% Xuan Research Paper Writing - NeurIPS 2026 Template
%
% 使用说明:
% 1. 复制此模板开始写作
% 2. 按章节顺序调用 Skill 生成内容
% 3. 将生成的内容粘贴到对应位置
% 4. 最终使用 Paper Review Skill 检查一致性
%
% NeurIPS 2026 轨道选择（提交前取消对应行的注释）:
% - \usepackage{neurips_2026}                 % 默认: Main Track (双盲审稿)
% - \usepackage[main]{neurips_2026}           % Main Track
% - \usepackage[position]{neurips_2026}      % Position Paper Track
% - \usepackage[eandd]{neurips_2026}         % Evaluations & Datasets Track
% - \usepackage[creativeai]{neurips_2026}    % Creative AI Track
% - \usepackage[preprint]{neurips_2026}      % arXiv预印本 (默认)
% - \usepackage[main,final]{neurips_2026}    % 最终版 (录用后)
%%%%%%%%%%%%%%%%%%%%%%%%%%%%%%%%%%%%%%%%%%%%%%%%%%%%%%%%%%%%%%%%%%%%%%%%%%%%%%%

\documentclass{article}

%===============================================================================
% NeurIPS 2026 样式 - 选择你的轨道
%===============================================================================
% 默认: preprint 模式 (用于arXiv)
\usepackage[preprint]{neurips_2026}

% 提交时取消下面一行的注释 (Main Track, 双盲审稿):
% \usepackage{neurips_2026}

% 其他轨道 (根据需要选择):
% \usepackage[position]{neurips_2026}      % Position Paper
% \usepackage[eandd]{neurips_2026}         % Evaluations & Datasets
% \usepackage[creativeai]{neurips_2026}    % Creative AI

% Workshop (如果需要):
% \usepackage[dblblindworkshop]{neurips_2026}
% \workshoptitle{Workshop Name}

%===============================================================================
% 推荐宏包
%===============================================================================
\usepackage[utf8]{inputenc} % allow utf-8 input
\usepackage[T1]{fontenc}    % use 8-bit T1 fonts
\usepackage{hyperref}       % hyperlinks
\usepackage{url}            % simple URL typesetting
\usepackage{booktabs}       % professional-quality tables
\usepackage{amsfonts}       % blackboard math symbols
\usepackage{amsmath}        % math environments
\usepackage{amssymb}        % math symbols
\usepackage{nicefrac}       % compact symbols for 1/2, etc.
\usepackage{microtype}      % microtypography
\usepackage{xcolor}         % colors
\usepackage{graphicx}       % figures
\usepackage{algorithm}      % algorithm environment
\usepackage{algpseudocode}  % algorithm pseudocode
\usepackage{caption}        % caption customization
\usepackage{subcaption}     % subfigures
\usepackage{multirow}       % multi-row tables
\usepackage{array}          % table column types
\usepackage{bm}             % bold math symbols

% 解决 hyperref 与 natbib 的冲突
\usepackage[capitalize,noabbrev]{cleveref}

%===============================================================================
% 图表路径
%===============================================================================
\graphicspath{{./figures/}{./images/}}

%===============================================================================
% 常用命令定义
%===============================================================================
\newcommand{\ours}{Ours}           % 替换为你的方法名称
\newcommand{\ie}{\textit{i.e.}}     % 即
\newcommand{\eg}{\textit{e.g.}}     % 例如
\newcommand{\etal}{\textit{et al.}} % 等人
\newcommand{\aka}{{a.k.a.}}         % 又称

% 数学简写 (按需添加)
\newcommand{\reals}{\mathbb{R}}
\newcommand{\naturals}{\mathbb{N}}
\newcommand{\integers}{\mathbb{Z}}

%===============================================================================
% 论文信息
%===============================================================================
\title{Your Paper Title Here}

% 作者信息
% 注意: 提交时使用匿名版本 (NeurIPS会自动处理)
% 最终版使用真实信息
\author{%
  First Author\thanks{Equal contribution} \quad
  Second Author$^*$ \quad
  Third Author \\[2pt]
  \textbf{Fourth Author \quad Fifth Author} \\[2pt]
  Affiliation \\[2pt]
  City, Country \\[2pt]
  \texttt{\{email1, email2\}@affiliation.com}
}

\begin{document}

\maketitle

%===============================================================================
% Abstract
% Skill: Abstract Skill (建议最后写)
% 输入: 全文完成后，提取核心贡献
% 注意: NeurIPS摘要限制一段，缩进0.5英寸，11pt行距
%===============================================================================
\begin{abstract}
% [SKILL OUTPUT: Abstract]
% 提示词: "基于我的 Introduction 和 Method，帮我写 Abstract"
%
% 结构 (NeurIPS格式):
% - 一句话概括: 任务 + 挑战
% - 核心方法: 我们的方法 + 关键技术 (1-2句)
% - 主要结果: 在XX数据集上提升XX% (1句)
% - 可选: 代码/数据可用性 (1句)
%
% 长度: 约150-250词

\textbf{TODO}: Call Abstract Skill after completing other sections.

% 示例 (替换为实际内容):
% We address the challenge of [task] by proposing [method], 
% which [key innovation]. Unlike previous approaches that 
% [limitation], our method [technical advantage]. Experiments 
% on [datasets] demonstrate [main results], outperforming 
% state-of-the-art methods by [margin] on [metric]. 
% Code is available at \url{https://github.com/xxx/xxx}.

\end{abstract}

%===============================================================================
% Introduction
% Skill: Introduction Skill
% 输入: 背景 + 主要内容 + 贡献 (自由文本)
% 输出: 5段式 Introduction
%===============================================================================
\section{Introduction}
\label{sec:intro}

% [SKILL OUTPUT: Introduction - Paragraph 1: Task and Applications]
% 提示词: "帮我写 Introduction。背景: [粘贴背景]
%          主要内容: [粘贴内容]
%          贡献: [粘贴贡献]"

\textbf{TODO}: Call Introduction Skill with background, content, and contributions.

% 示例结构 (根据实际内容调整):
% Para 1: Task definition + Applications (why important)
% Para 2: Existing methods + Their limitations (technical challenge)
% Para 3: Our method overview (high-level approach)
% Para 4: Technical contributions + Advantages
% Para 5: Experimental validation + Summary of results

%===============================================================================
% Related Work
% Skill: Related Work Skill
% 输入: 核心文献(3-5篇) + 候选文献池
% 依赖: Paper Summarizer Skill (预处理候选文献)
% 输出: 按技术主题组织的 Related Work
%===============================================================================
\section{Related Work}
\label{sec:related}

% [SKILL OUTPUT: Related Work]
% 提示词: "帮我写 Related Work
%          核心文献: [提供3-5篇最关键PDF]
%          候选文献: [提供文件夹或列表]"
%
% NeurIPS格式: 使用 \textbf{Topic Name}. 作为段落开头
%
% 结构:
% \textbf{Topic A Name.} [2-3 sentences for core papers, 1 for related]
% \textbf{Topic B Name.} [Similar structure]
% \textbf{Relation to Our Work.} [Optional summary]

\textbf{TODO}: Call Related Work Skill with core and candidate papers.

%===============================================================================
% Method
% Skill: Method Skill
% 输入: 代码仓库/文件 + Introduction 上下文
% 输出: 完整 Method 章节
%===============================================================================
\section{Method}
\label{sec:method}

% [SKILL OUTPUT: Method - Overview]
% 提示词: "帮我写 Method
%          代码路径: [提供代码文件夹]
%          Introduction 中的方法描述: [粘贴相关段落]"

\textbf{TODO}: Call Method Skill with code repository and Introduction context.

% NeurIPS推荐结构 (根据方法复杂度调整):
%
% \subsection{Overview}
% \label{sec:overview}
% [Pipeline summary + Figure reference]
%
% \subsection{[Module A Name]}
% \label{sec:module_a}
% \textbf{Motivation.} [Why this module?]
% \textbf{Design.} [How it works + formulation]
% \textbf{Advantages.} [Why better than alternatives]
%
% \subsection{[Module B Name]}
% \label{sec:module_b}
% [Same structure: Motivation → Design → Advantages]
%
% \subsection{Training Objective}
% \label{sec:objective}
% [Loss functions, optimization]
%
% \subsection{Implementation Details}
% \label{sec:implementation}
% [Architecture, hyperparameters, reproducibility info]

%===============================================================================
% Experiments
% Skill: Experiments Skill
% 输入: 实验草稿 + baseline论文 + 实验数据
% 阶段: B(设计) → A(填充) → C(完善)
%===============================================================================
\section{Experiments}
\label{sec:exp}

% Phase B: 实验设计阶段
% 提示词: "帮我设计 Experiments 章节
%          实验计划: [粘贴草稿]
%          Baseline论文: [提供PDF用于重构]"

\subsection{Experimental Setup}
\label{sec:setup}

\textbf{Datasets.}
% [SKILL OUTPUT: Dataset descriptions with statistics]
\textbf{Evaluation Metrics.}
% [SKILL OUTPUT: Metric definitions with formulas]
\textbf{Baselines.}
% [SKILL OUTPUT: Baseline descriptions, rewritten from their papers]
\textbf{Implementation Details.}
% [SKILL OUTPUT: Hyperparameters, architectures, training config]

\subsection{Main Results}
\label{sec:main}

% Phase A: 数据填充阶段
% 提示词: "填充实验数据
%          数据文件: [提供CSV/Excel/JSON]"
%
% NeurIPS表格格式:
% \begin{table}[t]
% \centering
% \caption{Caption should explain what the table shows and key takeaway.}
% \label{tab:main}
% \begin{tabular}{lccc}
% \toprule
% Method & Metric$_1$($\uparrow$) & Metric$_2$($\downarrow$) & Params \\
% \midrule
% Baseline A & 0.85 & 0.12 & 10M \\
% Baseline B & 0.87 & 0.10 & 15M \\
% \midrule
% Ours & \textbf{0.92} & \textbf{0.07} & 12M \\
% \bottomrule
% \end{tabular}
% \end{table}

\textbf{TODO}: Call Experiments Skill with data files.

\subsection{Ablation Study}
\label{sec:ablation}

% [SKILL OUTPUT: Ablation table + Analysis paragraphs]

\subsection{Additional Analysis}
\label{sec:analysis}

% [Optional: Efficiency, scalability, qualitative results, etc.]

%===============================================================================
% Conclusion
% Skill: Conclusion Skill
% 输入: 全文主要贡献
%===============================================================================
\section{Conclusion}
\label{sec:conclusion}

% [SKILL OUTPUT: Conclusion]
% 提示词: "帮我写 Conclusion，基于全文贡献"
%
% 结构 (1-2段):
% 1. 总结主要贡献 (呼应Introduction)
% 2. 主要实验发现
% 3. 局限与未来工作 (可选，NeurIPS推荐包含)

\textbf{TODO}: Call Conclusion Skill after completing other sections.

%===============================================================================
% Acknowledgments (最终版使用，提交时自动隐藏)
%===============================================================================
\begin{ack}
% 提交时此部分会自动隐藏 (使用 neurips_2026 样式时)
% 最终版使用真实致谢信息:
% This work was supported by ...
% We thank ... for helpful discussions.
%
% 需要披露资助信息:
% More info: https://neurips.cc/Conferences/2026/PaperInformation/FundingDisclosure
\end{ack}

%===============================================================================
% References
% NeurIPS: 使用小字体，无编号的一级标题
%===============================================================================
\section*{References}

% 参考文献用小字体 (9pt)
{
\small

% [SKILL OUTPUT: 自动生成或从你的 .bib 文件编译]

% 示例格式 (使用 natbib):
% \bibliographystyle{plainnat}
% \bibliography{references}

% 或手动列出:
% [1] Author, A., Author, B., & Author, C. (Year). Title. \textit{Venue}, pages.

}

%===============================================================================
% Appendix (可选，NeurIPS不限制附录页数)
% 注意: 论文主体必须能独立存在，关键实验不应只放在附录
%===============================================================================
\appendix

\section{Additional Experimental Results}
\label{sec:appendix-results}

% 补充实验结果，如:
% - 更多数据集上的结果
% - 超参数敏感性分析
% - 额外的消融实验

\section{Implementation Details}
\label{sec:appendix-implementation}

% 更多实现细节，如:
% - 网络架构详细配置
% - 训练曲线
% - 伪代码细节

\section{Proofs}
\label{sec:appendix-proofs}

% 定理证明 (如果有)

%===============================================================================
% NeurIPS Checklist (提交时必需)
% 复制 checklist.tex 的内容或 \input{checklist}
%===============================================================================
% \newpage
% \input{checklist}

\end{document}
